\section*{Resumen.}

En el presente trabajo se describe el diseño, desarrollo e implementación del Sistema de Expedientes de Trámites y Citas (SETC), una plataforma web empresarial e institucional concebida para modernizar, digitalizar y transparentar los procesos de radicación, evaluación técnica documental y resolución de trámites ciudadanos, así como la gestión integral de citas presenciales y virtuales. El sistema está desarrollado bajo el patrón arquitectónico Modelo-Vista-Controlador (MVC) utilizando el framework Laravel 13 sobre PHP 8.4, complementado con Tailwind CSS y Alpine.js para la construcción de una interfaz reactiva, intuitiva y responsiva, y un modelo de datos relacional normalizado.

La solución erradica las aglomeraciones en ventanillas físicas, elimina el consumo excesivo de papel y reduce sustancialmente los tiempos de espera a través de un flujo estructurado en cinco fases procedimentales sincronizadas en vivo: radicación telemática con validación estricta de recaudos digitales en formato PDF (hasta 10 MB) y emisión atómica de folios institucionales (\texttt{SETC-YYYY-XXXXX}), asignación técnica pericial, mesa de dictamen con calificación granular de requisitos y gestión reactiva de subsanaciones, formalización de dictamen técnico y expedición de resolución oficial en formato PDF provista de código QR de verificación pública. Asimismo, la plataforma incorpora un Asistente de Inteligencia Artificial (IA) interactivo diseñado para orientar de manera continua al ciudadano, absolver dudas sobre los recaudos normativos y brindar soporte contextual durante el llenado de su solicitud. Para los niveles de supervisión y administración, la plataforma dispone de una matriz analítica de distribución de trabajo, garantizando la rendición de cuentas mediante un control de acceso estricto basado en roles (RBAC) y una bitácora transaccional forense inmutable. La robustez y seguridad del sistema fueron certificadas mediante una suite de 72 pruebas automatizadas en Pest PHP v3 y un despliegue contenerizado en Docker.
